\section{Sécurité appliquée}

\subsection{Authentification de l'application web}

L'application React est un client propriétaire de la même plateforme. Elle utilise
donc Laravel Sanctum en mode \textbf{session stateful} plutôt qu'un JWT conservé dans
le stockage du navigateur. Avant une connexion locale, le frontend récupère le
cookie CSRF et envoie les requêtes avec les cookies activés. Le cookie de session est
\texttt{HttpOnly}, ce qui réduit l'exposition du secret de session au code
JavaScript.

Google OAuth est orchestré par le backend. Le secret client Google n'est jamais
fourni au frontend. Après le callback, Laravel exige un email vérifié, lie ou crée le
compte local puis établit la même session Sanctum. Un compte Google ne reçoit aucun
rôle privilégié du fournisseur : les rôles et l'organisation sont attribués
uniquement par la plateforme. Dans le périmètre actuel, Microsoft OAuth est différé.

Les mots de passe locaux sont hashés par le mécanisme Laravel. Les formulaires de
récupération répondent de manière générique afin de limiter l'énumération des
comptes. Les comptes exclusivement Google n'affichent pas une action de changement
de mot de passe tant qu'une authentification locale n'a pas été activée.

\subsection{Autorisation et périmètre organisationnel}

L'autorisation repose sur trois niveaux complémentaires :

\begin{enumerate}[leftmargin=8mm]
  \item les permissions Spatie décrivent les capacités d'un rôle ;
  \item les policies contrôlent l'action sur la ressource concrète ;
  \item \texttt{canAccessOrganization} et les filtres de requête garantissent le
  périmètre de données.
\end{enumerate}

\begin{table}[H]
  \centering
  \caption{Périmètre d'accès des rôles}
  \label{tab:role-data-scope}
  \begin{tabularx}{\textwidth}{|p{2.8cm}|p{3.3cm}|Y|}
    \hline
    \rowcolor{ctmint}
    \textbf{Rôle} & \textbf{Association} & \textbf{Périmètre}\\\hline

    Super administrateur & Global, sans organisation &
    Organisations, intégrations, audit et commerce SaaS selon ses permissions.\\\hline
    Administrateur & Une organisation &
    Employés, clients, actifs, tarifs, commerce et rapports de son organisation.\\\hline
    Opérateur & Une organisation &
    Stations, connecteurs, sessions, alertes, maintenances et interventions de son
    organisation.\\\hline
    Technicien & Une organisation &
    Consultation des actifs et traitement des alertes, interventions ou maintenances
    autorisées, sans création de station.\\\hline
    Client & Global, sans organisation &
    Stations publiques et uniquement ses propres tentatives, sessions, paiements,
    abonnements et notifications.\\\hline

  \end{tabularx}
\end{table}

Le frontend adapte la navigation, mais l'absence d'un bouton ne constitue jamais une
autorisation. Chaque endpoint sensible appelle une policy. L'administrateur ne peut
pas gérer les employés d'une autre organisation, changer son propre rôle, désactiver
son propre compte ni supprimer le dernier administrateur actif. Les identifiants
présents dans l'URL sont donc revalidés afin de prévenir les accès directs de type
IDOR.

\subsection{Sécurité des frontières OCPP}

Deux niveaux protègent la communication des bornes :

\begin{itemize}[leftmargin=7mm]
  \item la borne utilise une identité OCPP et un secret Basic Auth. Pour une borne
  externe, le secret aléatoire est affiché une seule fois et seul son hash est
  stocké ;
  \item la passerelle signe ses requêtes internes avec HMAC-SHA256 sur le corps
  exact, un timestamp Unix et un identifiant UUID.
\end{itemize}

Laravel rejette les signatures invalides, les timestamps trop anciens et les UUID
déjà utilisés. L'identifiant métier de l'événement rend également l'ingestion
idempotente : une répétition identique retourne l'enregistrement existant, tandis
qu'un même identifiant associé à un contenu différent est refusé. Les payloads de
commandes sont chiffrés au repos et expirent. Un seul gateway peut réclamer une
commande en attente.

HMAC garantit l'intégrité et l'authenticité du service, mais pas la confidentialité.
En production, la borne utilise donc \texttt{wss://} et le gateway communique avec
Laravel par HTTPS ou au sein d'un réseau privé chiffré. Le secret partagé du
simulateur n'est acceptable que dans l'environnement local ; chaque borne physique
doit recevoir un secret individuel.

\subsection{Sécurité du paiement}

Le simulateur n'accepte aucune carte réelle. L'API externe reçoit une référence, un
montant, une devise, un résultat de simulation et une clé d'idempotence. Le webhook
retour est signé et limité en fréquence. \texttt{PaymentProviderEvent} conserve
l'identifiant du fournisseur, la signature vérifiée et le résultat de
réconciliation. Les règles suivantes s'appliquent :

\begin{itemize}[leftmargin=7mm]
  \item aucune confirmation du frontend ne marque directement un paiement comme
  réussi ;
  \item un webhook invalide n'est pas traité ;
  \item une répétition valide ne produit pas une seconde opération ;
  \item un délai ou une indisponibilité technique est distingué d'un refus métier ;
  \item les secrets, numéros de carte et payloads sensibles ne sont jamais journalisés.
\end{itemize}

L'intégration d'un prestataire réel nécessitera en plus la conformité aux exigences
du fournisseur, TLS public, une politique de remboursement, la réconciliation
financière et, selon le mode choisi, l'externalisation complète des données de carte.

\subsection{Validation, limitation et gestion des erreurs}

Les écritures utilisent des Form Requests et retournent une erreur métier contrôlée
plutôt qu'une trace SQL ou une exception interne. Des buckets de limitation
indépendants empêchent un parcours public de bloquer un autre :

\begin{table}[H]
  \centering
  \caption{Principales limites applicatives}
  \label{tab:rate-limits}
  \begin{tabularx}{\textwidth}{|Y|p{3.1cm}|Y|}
    \hline
    \rowcolor{ctmint}
    \textbf{Flux} & \textbf{Limite} & \textbf{Clé}\\\hline

    Demande de démonstration & 3 par heure & Adresse IP\\\hline
    Inspection d'invitation & 10 par minute & Adresse IP\\\hline
    Acceptation d'invitation & 5 par minute & IP et hash de l'email\\\hline
    Gestion des invitations employé & 10 par minute & Utilisateur authentifié\\\hline
    API interne OCPP & 600 par minute & Adresse de la passerelle\\\hline
    Webhook de paiement & 120 par minute & Adresse du fournisseur\\\hline
    Changement de mot de passe & 5 par minute & Session authentifiée\\\hline

  \end{tabularx}
\end{table}

La déduplication métier complète le throttling. Par exemple, une demande récente
ayant la même identité n'est pas recréée, et les alertes automatiques utilisent une
clé stable. Cette combinaison traite à la fois les abus volumétriques et les doubles
clics légitimes.

\subsection{Secrets, fichiers et journalisation}

Les fichiers \texttt{.env} contiennent les clés Laravel, Google, Resend, Reverb,
OCPP et paiement. Ils sont ignorés par Git ; seuls les \texttt{.env.example} sans
valeurs privées sont versionnés. Le frontend ne reçoit que les paramètres publics
préfixés \texttt{VITE\_}. Les secrets de bornes, tokens OAuth, mots de passe,
en-têtes \texttt{Authorization} et signatures ne doivent pas apparaître dans les
logs.

Les photos d'intervention, rapports et documents techniques utilisent le disque
Laravel privé. Un contrôleur autorisé lit le fichier et le diffuse après policy ;
une URL de stockage publique ne contourne donc pas le périmètre organisationnel. Les
types MIME, tailles et extensions sont validés avant enregistrement.

\texttt{PlatformAuditLog} conserve les actions sensibles : gestion de comptes et
rôles, changement de paramètres, commissioning, rotation de credentials, tarifs,
abonnements et états financiers. L'auteur, l'organisation, la cible, l'action et un
résumé sont enregistrés sans inclure le secret modifié. Une tâche quotidienne
applique la politique de rétention.

\subsection{Analyse synthétique des risques}

\begin{longtable}{|p{3.2cm}|p{4.4cm}|p{6.6cm}|}
  \caption{Risques principaux et mesures de réduction}
  \label{tab:security-risk-summary}\\
  \hline
  \textbf{Risque} & \textbf{Exemple} & \textbf{Mesures}\\
  \hline
  \endfirsthead
  \hline
  \textbf{Risque} & \textbf{Exemple} & \textbf{Mesures}\\
  \hline
  \endhead
  Vol de session & Script lisant un token du navigateur &
  Cookie HttpOnly, Sanctum stateful, CSRF, rotation et TLS.\\\hline
  IDOR multi-tenant & Accès à \texttt{/stations/11} d'une autre organisation &
  Policies, scoping organisationnel, tests d'isolation.\\\hline
  Rejeu OCPP & Réémission d'un événement signé &
  Timestamp, UUID de requête, HMAC et identifiant d'événement idempotent.\\\hline
  Faux paiement & Callback frontend ou webhook forgé &
  Confirmation serveur, signature HMAC et événements fournisseur uniques.\\\hline
  Double facturation & Répétition d'une capture après timeout &
  Clés d'idempotence et réconciliation persistée.\\\hline
  Exposition documentaire & URL publique d'une photo technique &
  Stockage privé, validation et téléchargement soumis à policy.\\\hline
  Secret dans Git & Commit d'un client secret OAuth &
  \texttt{.gitignore}, exemples vides et injection par environnement.\\\hline
  Déni de service logique & Soumission répétée de demandes ou invitations &
  Rate limiting indépendant, honeypot et déduplication.\\
  \hline
\end{longtable}

\subsection{Limites de sécurité du prototype}

Le socle protège les parcours implémentés, mais ne remplace pas un audit de
production. Il reste notamment à automatiser un scan des dépendances, renforcer les
en-têtes HTTP, exécuter des tests de charge OCPP, définir la rotation des secrets,
déployer une terminaison TLS, superviser les processus et tester les sauvegardes.
Une borne physique et un fournisseur de paiement réel devront être homologués
séparément. Ces limites sont documentées afin de ne pas confondre un MVP sécurisé
par conception avec une certification d'exploitation.
